% ---
% title: Do the two stocks move independently?
% tags: finance, independent-rv, joint-discrete
% ---

\begin{problem}
	\begin{context}
		Investors spread their money over several stocks to reduce risk. This
		works best when the stocks do not tend to move together. If two stocks
		usually rise and fall on the same days, holding both protects much less
		than holding two stocks that move independently of each other.
	\end{context}
	These are the two stocks of Exercise~\exref{two-stocks}. On a given day,
	let $X=1$ if stock A rises and $X=-1$ if it falls. Let $Y=1$, $0$ or $-1$
	if stock B rises, stays unchanged or falls. The joint probability mass
	function $p(x,y)=P(X=x,Y=y)$ is given by the table
	\begin{center}
		\renewcommand{\arraystretch}{1.2}
		\begin{tabular}{c|ccc}
			         & $Y=-1$ & $Y=0$  & $Y=1$  \\
			\hline
			$X=-1$   & $0.20$ & $0.10$ & $0.10$ \\
			$X=1$    & $0.05$ & $0.15$ & $0.40$
		\end{tabular}
	\end{center}
	(a) Are $X$ and $Y$ independent?
	(b) Write down the joint probability mass function that $X$ and $Y$ would
	have if they were independent, with the same marginal probability mass
	functions.
	(c) For these independent stocks, find $P(X=Y)$, the probability that both
	fall or both rise. Compare it with $P(X=Y)=0.60$ for the real stocks.
\end{problem}

\begin{slideproblem}
	Stocks A and B of Exercise~\exref{two-stocks}: $X=\pm1$ if A
	rises/falls, $Y=1,0,-1$ if B rises/stays/falls. Joint pmf with the
	marginals in the margins:
	\begin{center}
		\renewcommand{\arraystretch}{1.15}
		\begin{tabular}{c|ccc|c}
			       & $y=-1$ & $y=0$  & $y=1$  & $p_X(x)$ \\
			\hline
			$x=-1$ & $0.20$ & $0.10$ & $0.10$ & $0.4$ \\
			$x=1$  & $0.05$ & $0.15$ & $0.40$ & $0.6$ \\
			\hline
			$p_Y(y)$ & $0.25$ & $0.25$ & $0.5$ & $1$
		\end{tabular}
	\end{center}
	(a) Are $X$ and $Y$ independent?
	(b) Write down the joint pmf of independent $X$, $Y$ with the same
	marginals.
	(c) For those, find $P(X=Y)$ and compare it with $0.60$ for the real
	stocks.
\end{slideproblem}

\begin{slides}
	\small
	\exstatement
	\pause
	\sol (a) Independent $\iff$ $p(x,y)=p_X(x)\,p_Y(y)$ in \alert{every}
	cell. But\\
	$p(-1,-1)=0.20\neq0.10=0.4\cdot0.25=p_X(-1)\,p_Y(-1)$:
	\alert{not independent}.\\
	\pause
	One bad cell is enough for ``no''; for ``yes'' \emph{all} cells must match.
\end{slides}

\begin{slides}
	\small
	(b) Independent stocks, same marginals: each cell is
	\alert{row sum $\times$ column sum}.
	\begin{center}
		\renewcommand{\arraystretch}{1.25}
		\begin{tabular}{c|ccc|c}
			       & $y=-1$ & $y=0$  & $y=1$  & $p_X(x)$ \\
			\hline
			$x=-1$ & \uncover<2->{$0.10$} & \uncover<2->{$0.10$}
			& \uncover<2->{$0.20$} & $0.4$ \\
			$x=1$  & \uncover<2->{$0.15$} & \uncover<2->{$0.15$}
			& \uncover<2->{$0.30$} & $0.6$ \\
			\hline
			$p_Y(y)$ & $0.25$ & $0.25$ & $0.5$ & $1$
		\end{tabular}
	\end{center}
	\uncover<2->{For example $0.4\cdot0.5=0.20$ and $0.6\cdot0.25=0.15$.}\\[1mm]
	\pause\pause
	(c) Both fall or both rise:
	\[ P(X=Y)=p(-1,-1)+p(1,1)=0.10+0.30=\alert{0.40}\;<\;0.60. \]
	\pause
	The real stocks move in the same direction more often than independent
	ones would: they tend to \alert{move together}.
\end{slides}

\begin{solution}
	\idea Two discrete random variables are independent exactly when every
	cell of the joint table equals the product of its row sum and column sum:
	$p(x,y)=p_X(x)\,p_Y(y)$ for all $x$ and $y$. To show dependence, one cell
	where this fails is enough. To build an independent pair with given
	marginals, fill every cell with this product.

	\step{(a)} First we need the marginals. They are the row sums and the
	column sums of the table (as in Exercise~\exref{two-stocks}):
	\[ p_X(-1)=0.20+0.10+0.10=0.4,\qquad p_X(1)=0.05+0.15+0.40=0.6, \]
	\[ p_Y(-1)=0.20+0.05=0.25,\qquad p_Y(0)=0.10+0.15=0.25,\qquad
		p_Y(1)=0.10+0.40=0.5. \]
	Now compare a cell with the product of its marginals. Take the cell where
	both stocks fall:
	\[ p(-1,-1)=0.20,\qquad\text{but}\qquad
		p_X(-1)\,p_Y(-1)=0.4\cdot0.25=0.10. \]
	These are different, so $P(X=-1,\,Y=-1)\neq P(X=-1)\,P(Y=-1)$, and $X$
	and $Y$ are not independent. The two stocks fall together twice as often
	as independent stocks with the same marginals would.

	Note the asymmetry of the check. To prove that $X$ and $Y$ are
	\emph{dependent}, a single cell that fails is enough. To prove that they
	are \emph{independent}, we would have to check that \emph{all} six cells
	match.

	\step{(b)} If $X$ and $Y$ were independent with the same marginals, every
	cell would be the product of its row sum and its column sum:
	\[ \begin{aligned}
		p(-1,-1)&=0.4\cdot0.25=0.10, & p(-1,0)&=0.4\cdot0.25=0.10, &
		p(-1,1)&=0.4\cdot0.5=0.20,\\
		p(1,-1)&=0.6\cdot0.25=0.15, & p(1,0)&=0.6\cdot0.25=0.15, &
		p(1,1)&=0.6\cdot0.5=0.30.
	\end{aligned} \]
	In a table, with the marginals in the margins:
	\begin{center}
		\renewcommand{\arraystretch}{1.25}
		\begin{tabular}{c|ccc|c}
			         & $y=-1$ & $y=0$  & $y=1$  & $p_X(x)$ \\
			\hline
			$x=-1$   & $0.10$ & $0.10$ & $0.20$ & $0.4$ \\
			$x=1$    & $0.15$ & $0.15$ & $0.30$ & $0.6$ \\
			\hline
			$p_Y(y)$ & $0.25$ & $0.25$ & $0.5$  & $1$
		\end{tabular}
	\end{center}
	Check: the rows still add up to $0.4$ and $0.6$, and the columns to
	$0.25$, $0.25$ and $0.5$. So this is a valid joint pmf with the same
	marginals, but a different joint distribution. The marginals alone do not
	determine the joint distribution.

	\step{(c)} $X=Y$ happens when both stocks fall, $(-1,-1)$, or both rise,
	$(1,1)$; the value $Y=0$ never equals $X$. These two cells are mutually
	exclusive, so we add them, now in the independent table:
	\[ P(X=Y)=p(-1,-1)+p(1,1)=0.10+0.30=0.40. \]
	For the real stocks the same two cells give $0.20+0.40=0.60$. So the real
	stocks move in the same direction on 60\% of the days, but independent
	stocks with the same marginals would do so on only 40\% of the days. The
	real stocks tend to move together.
	\begin{answer}
		(a) No: $p(-1,-1)=0.20\neq0.10=p_X(-1)\,p_Y(-1)$
		\quad (b) $x=-1$: $0.10,\ 0.10,\ 0.20$; \ $x=1$: $0.15,\ 0.15,\ 0.30$
		(for $y=-1,0,1$)
		\quad (c) $P(X=Y)=0.40$, less than $0.60$ for the real stocks
	\end{answer}
	\begin{remark}
		\textbf{Common mistake:} checking only one cell and concluding
		independence because it happens to match. For example, the cell
		$(-1,0)$ gives $p(-1,0)=0.10=0.4\cdot0.25$, yet $X$ and $Y$ are not
		independent. A match in one cell proves nothing; independence needs
		$p(x,y)=p_X(x)\,p_Y(y)$ in every cell.
	\end{remark}
\end{solution}
